\begin{ouverture}[geometrie]{16}{Translations, agrandissements et réductions}{Partie 4 · Géométrie et espace}
{Glissement · Frises · Longueurs · Aires · Volumes}
\objectif{Construire l’image d’une figure par translation.}
\objectif{Utiliser les propriétés de conservation.}
\objectif{Relier un rapport de longueurs aux aires et aux volumes.}
\prerequis{Repérage, parallélogrammes et proportionnalité.}
\end{ouverture}
\pagestyle{otMain}\essentiel\fontsize{11}{15}\selectfont
\phantomsection\label{cours-16-1}
\essSection{1}{Faire glisser une figure}
\begin{definition}Une translation fait glisser une figure selon une même direction, un même sens et une même longueur, sans la faire tourner. La translation qui transforme P en Q applique ce même déplacement à chaque point de la figure.\end{definition}
\label{fig-cours-16-1-cours-translation}\figureGeo{cours-translation}

\begin{propriete}Une translation conserve les longueurs, les angles, l’alignement, le parallélisme et les aires. Un segment et son image sont parallèles et de même longueur. Si A a pour image A’ et B pour image B’, le quadrilatère A-B-B’-A’ est un parallélogramme quand ses sommets sont distincts et non alignés.\end{propriete}
\begin{savoirfaire}[Construire]Sur quadrillage, compter le déplacement de P vers Q et le reproduire depuis chaque sommet. Sans quadrillage, construire les parallèles et reporter la même longueur dans le même sens. Une frise ou un pavage peut répéter un motif par translations.\end{savoirfaire}
\clearpage
\phantomsection\label{cours-16-2}
\essSection{2}{Changer de taille}
\begin{definition}Un agrandissement ou une réduction multiplie toutes les longueurs correspondantes par un même rapport positif $k$. Si $k>1$, c’est un agrandissement ; si $0<k<1$, une réduction. Les angles sont conservés.\end{definition}
\begin{propriete}Quand les longueurs sont multipliées par $k$, les périmètres sont multipliés par $k$, les aires par $k^2$ et les volumes par $k^3$. On peut retrouver ces facteurs en appliquant les formules d’aire et de volume aux nouvelles dimensions.\end{propriete}
\begin{exemple}Pour un rectangle de 6 cm sur 4 cm agrandi avec $k=1{,}5$, les côtés deviennent 9 et 6 cm. L’aire passe de 24 à 54 cm$^2$ : elle est multipliée par $1{,}5^2=2{,}25$. Pour un cube réduit de moitié, le volume est divisé par 8.\end{exemple}
\begin{attention}Une translation conserve la taille, contrairement à un agrandissement. Pour retrouver $k$ depuis les aires, on cherche le nombre positif dont le carré est le rapport des aires ; le rapport des aires n’est pas directement le rapport des longueurs.\end{attention}

\clearpage\rubrique{Découverte}
\begin{blocLivre}{ex-16-1}\begin{exercice}[1]{Décrire le glissement}
Dans le quadrillage, la translation transforme P(0;0) en Q(3;1). Quelle est l’image de A(1;2) ?
\label{fig-ex-16-1}\figureGeo{trans-01}
\end{exercice}\end{blocLivre}
\begin{blocLivre}{ex-16-2}\begin{exercice}[1]{Construire un triangle image}
Construire l’image de $A(1;1)$, $B(3;1)$, $C(2;3)$ par le glissement de 3 carreaux à droite et 2 vers le haut.
\label{fig-ex-16-2}\figureGeo{trans-02}
\end{exercice}\end{blocLivre}
\begin{blocLivre}{ex-16-3}\begin{exercice}[1]{Un déplacement négatif}
Une translation déplace de 4 carreaux vers la gauche et 1 vers le bas. Donner l’image de B(2;3).
\label{fig-ex-16-3}\figureGeo{trans-03}
\end{exercice}\end{blocLivre}
\begin{blocLivre}{ex-16-4}\begin{exercice}[1]{Conserver des mesures}
ABC a $AB=5$ cm, $\widehat{BAC}=42^\circ$ et une aire de 9 cm$^2$. Donner ces mesures dans son image A’B’C’ par translation.
\label{fig-ex-16-4}\figureGeo{trans-04}
\end{exercice}\end{blocLivre}
\begin{blocLivre}{ex-16-5}\begin{exercice}[1]{Reconnaître une frise}
Les débuts de cinq motifs identiques d’une frise sont espacés de 4 cm. Quel glissement mène du premier au suivant ? Du premier au cinquième ?
\label{fig-ex-16-5}\figureGeo{trans-05}
\end{exercice}\end{blocLivre}

\rubrique{Entraînement}
\begin{blocLivre}{ex-16-6}\begin{exercice}[2]{Un parallélogramme}
A a pour image A’ et B pour image B’ par une même translation. Les quatre points sont non alignés et distincts. Justifier que A-B-B’-A’ est un parallélogramme.
\label{fig-ex-16-6}\figureGeo{trans-06}
\end{exercice}\end{blocLivre}
\begin{blocLivre}{ex-16-7}\begin{exercice}[2]{Agrandir une figure}
Un rectangle de 6 cm sur 4 cm est agrandi avec un rapport de 1,5. Donner ses dimensions et son périmètre.
\label{fig-ex-16-7}\figureGeo{trans-07}
\end{exercice}\end{blocLivre}
\begin{blocLivre}{ex-16-8}\begin{exercice}[2]{Réduire un triangle}
Un triangle de côtés 5, 7 et 8 cm est réduit au rapport $2/5$. Donner les trois nouvelles longueurs.
\label{fig-ex-16-8}\figureGeo{trans-08}
\end{exercice}\end{blocLivre}
\begin{blocLivre}{ex-16-9}\begin{exercice}[2]{Une aire agrandie}
Une figure d’aire 18 cm$^2$ est agrandie au rapport 3. Quelle est sa nouvelle aire ?
\label{fig-ex-16-9}\figureGeo{trans-09}
\end{exercice}\end{blocLivre}
\begin{blocLivre}{ex-16-10}\begin{exercice}[2]{Un volume réduit}
Un pavé droit de volume 240 cm$^3$ est réduit au rapport $1/2$. Quel est le volume image ?
\label{fig-ex-16-10}\figureGeo{trans-10}
\end{exercice}\end{blocLivre}

\rubrique{Problèmes}
\begin{blocLivre}{ex-16-11}\begin{exercice}[3]{Retrouver le rapport}
Deux figures sont l’une un agrandissement de l’autre. Leurs aires sont 20 et 45 cm$^2$. Trouver le rapport d’agrandissement et l’image d’un côté de 8 cm.
\label{fig-ex-16-11}\figureGeo{trans-11}
\end{exercice}\end{blocLivre}
\begin{blocLivre}{ex-16-12}\begin{exercice}[3]{Un pavé à l’échelle}
Un pavé droit mesure 12, 8 et 6 cm. On réduit toutes ses dimensions au tiers. Donner dimensions, aire de la face 12 par 8 et volume après réduction.
\label{fig-ex-16-12}\figureGeo{trans-12}
\end{exercice}\end{blocLivre}
\begin{blocLivre}{ex-16-13}\begin{exercice}[3]{Deux déplacements}
On déplace un point de 5 carreaux à droite et 2 en haut, puis de 3 à gauche et 4 en haut. Quel glissement unique donne le même résultat ?
\label{fig-ex-16-13}\figureGeo{trans-13}
\end{exercice}\end{blocLivre}
\begin{blocLivre}{ex-16-14}\begin{exercice}[3]{Le bon coefficient de peinture}
Une maquette de surface extérieure 0,8 m$^2$ est agrandie en multipliant toutes ses longueurs par 2,5. La peinture couvre 4 m$^2$ par litre. Quelle quantité théorique faut-il pour la nouvelle surface ?
\label{fig-ex-16-14}\figureGeo{trans-14}
\end{exercice}\end{blocLivre}
\begin{blocLivre}{ex-16-15}\begin{exercice}[3]{Un pavage par translations}
Un motif carré de côté 2 cm est répété sans trou ni chevauchement sur 6 colonnes et 4 lignes. Combien de motifs utilise-t-on ? Donner l’aire totale et un déplacement permettant de passer d’un motif au motif situé deux colonnes à droite et une ligne au-dessus.
\label{fig-ex-16-15}\figureGeo{trans-15}
\end{exercice}\end{blocLivre}

\clearpage
